\chapter{Fitch Cheney's Card Trick}

\section{The Story}
You are the magician. A volunteer shuffles a normal deck. Your partner stands
next to the table. You leave the room for a moment.

When you come back, you glance at four face-up cards and immediately name the
fifth card that is hidden face down. The audience freezes for a second, then
starts laughing, shouting, and accusing you of witchcraft.

But there is no mirror, no marked deck, and no secret talking. The miracle is
pure structure: a tiny encoding system hidden in card order.

For clarity, we will still call your partner \textbf{Alice} and the magician
\textbf{Bob}.

\section{What the Audience Sees}
\begin{enumerate}
  \item A volunteer shuffles the deck and randomly chooses \(5\) cards.
  \item Bob leaves the room, so he cannot see anything.
  \item Alice looks at the \(5\) chosen cards, hides one card face down, and
  lays the other \(4\) cards face up in a specific order.
  \item Bob comes back, looks only at those \(4\) visible cards, and names the
  hidden fifth card exactly.
\end{enumerate}

To the audience, it looks impossible. There is no marked deck, no mirrors, and
no code words. The secret is pure combinatorics.

\section{Reveal: How the Trick Works}
\subsection*{Key idea}
With \(4\) visible cards, Alice can arrange them in \(4! = 24\) different ways.
So she can encode one of \(24\) possibilities by order alone.

\subsection*{Step 1: Find two cards of the same suit}
Among \(5\) cards and only \(4\) suits, at least two cards share a suit
(pigeonhole principle). Call those two same-suit cards \(X\) and \(Y\).

One of them will be shown first (the \emph{key card}), and the other one will
be hidden (the \emph{target card}).

\subsection*{Step 2: Use circular rank distance}
Within one suit, order ranks in a cycle:
\[
2,3,\dots,10,J,Q,K,A.
\]
From the key card's rank to the target card's rank, count forward distance
\(1\) to \(12\). For any two distinct ranks, one direction is at most \(6\).

Alice chooses which card is key and which is target so the forward distance is
in \(\{1,2,3,4,5,6\}\). She places the key card first.

\subsection*{Step 3: Encode distances \(1\) to \(6\) with \(3!\) orders}
After fixing the first visible card, three cards remain. Their order can be one
of \(3!=6\) permutations, perfect for encoding distances \(1\) through \(6\).

Use this fixed comparison rule:
\begin{itemize}
  \item Compare rank first: \(2<3<\cdots<K<A\).
  \item If ranks are equal, compare suit:
  \(\clubsuit<\diamondsuit<\heartsuit<\spadesuit\).
\end{itemize}
Sort the three non-key cards from small to large by this rule, and call them
\(a<b<c\).

Then agree on this encoding table:
\begin{center}
\begin{tabular}{c l}
\hline
\textbf{Distance} & \textbf{Order of last three cards} \\
\hline
1 & \(a,b,c\) \\
2 & \(a,c,b\) \\
3 & \(b,a,c\) \\
4 & \(b,c,a\) \\
5 & \(c,a,b\) \\
6 & \(c,b,a\) \\
\hline
\end{tabular}
\end{center}

Bob sees:
\begin{itemize}
  \item the first card \(\Rightarrow\) suit of the hidden card,
  \item the order of the last three \(\Rightarrow\) distance \(1\) to \(6\).
\end{itemize}
So he reconstructs the exact hidden card.

\section{Mini Example}
Suppose the five selected cards are
\[
9\heartsuit,\ K\heartsuit,\ 3\spadesuit,\ J\diamondsuit,\ 5\clubsuit.
\]
Alice decides:
\begin{itemize}
  \item key card \(=9\heartsuit\) (shown first),
  \item hidden card \(=K\heartsuit\),
  \item remaining three cards: \(3\spadesuit,\ J\diamondsuit,\ 5\clubsuit\).
\end{itemize}

From \(9\) to \(K\), the forward distance is \(4\)
\((9\to10\to J\to Q\to K)\), so Alice must encode \(4\).

Using the fixed rule above:
\[
3\spadesuit < 5\clubsuit < J\diamondsuit,
\]
so
\[
a=3\spadesuit,\quad b=5\clubsuit,\quad c=J\diamondsuit.
\]
Distance \(4\) corresponds to order \(b,c,a\), so the four visible cards are
\[
9\heartsuit,\ 5\clubsuit,\ J\diamondsuit,\ 3\spadesuit.
\]
This means Bob sees the cards in this exact left-to-right order:
\[
\boxed{9\heartsuit}\ ,\ \boxed{5\clubsuit}\ ,\ \boxed{J\diamondsuit}\ ,\ \boxed{3\spadesuit}.
\]
From the first card, Bob gets the hidden suit \(\heartsuit\). From
\(5\clubsuit,J\diamondsuit,3\spadesuit\), he reads order \(b,c,a\), so distance
\(=4\). Therefore Bob announces \(K\heartsuit\).

\subsection*{Another Example (Hide the Lower Card)}
Now consider the same-suit pair \(2\clubsuit\) and \(K\clubsuit\). If Alice used
\(2\clubsuit\) as key and tried to hide \(K\clubsuit\), the forward distance
would be \(11\), which is too large for our \(1\) to \(6\) code.

Take a full hand:
\[
2\clubsuit,\ K\clubsuit,\ 7\heartsuit,\ 4\spadesuit,\ Q\diamondsuit.
\]
So Alice reverses the direction: she \emph{shows} \(K\clubsuit\) first and
\emph{hides} \(2\clubsuit\). Then the forward distance is
\[
K\to A\to 2,
\]
which is distance \(2\).

The remaining three are \(7\heartsuit,\ 4\spadesuit,\ Q\diamondsuit\). Under
the same fixed rule:
\[
4\spadesuit < 7\heartsuit < Q\diamondsuit,
\]
so
\[
a=4\spadesuit,\quad b=7\heartsuit,\quad c=Q\diamondsuit.
\]
Distance \(2\) corresponds to order \(a,c,b\), so Alice lays out:
\[
K\clubsuit,\ 4\spadesuit,\ Q\diamondsuit,\ 7\heartsuit
\]
with \(2\clubsuit\) hidden.
So Bob sees, in exact order:
\[
\boxed{K\clubsuit}\ ,\ \boxed{4\spadesuit}\ ,\ \boxed{Q\diamondsuit}\ ,\ \boxed{7\heartsuit}.
\]
From the first card he gets suit \(\clubsuit\). From
\(4\spadesuit,Q\diamondsuit,7\heartsuit\), he reads \(a,c,b\), so distance
\(=2\). Starting at \(K\clubsuit\), two steps forward in the cycle gives
\(2\clubsuit\), which is the hidden card.

This example is important: sometimes Alice must hide the numerically lower card
to keep the distance in \(\{1,2,3,4,5,6\}\).

\section{Perform It at Home}
You can learn this in one evening. Practice with a friend:
\begin{enumerate}
  \item Memorize the distance table \(1\) to \(6\).
  \item Agree on one strict ordering rule for comparing cards.
  \item Drill 20 random hands of 5 cards and time yourselves.
  \item Start slow and accurate, then speed up.
\end{enumerate}

After a little practice, you can perform this at dinner or a party. It is a
great example of how mathematics can feel like mind reading.

\section*{Reference}
The explanation follows the MIT note by Kleber \cite{kleber:fitch-cheney}.
